\chapter{Syarat-syarat Istinja dengan Batu}

\begin{matanbox}[Matan]
Syarat-syarat istinja dengan batu ada 8:

\begin{enumerate}
    \item Dengan tiga batu.
    \item Mampu membersihkan tempat keluarnya najis.
    \item Najis belum mengering. 
    \item Najis tidak berpindah tempat.
    \item Najis tidak mengenai dengan najis yang lain.
    \item Najis tidak keluar melebihi \textit{shafhah} dan \textit{hasyafah}. 
    \item Najis tidak terkena air.
    \item Menggunakan batu yang suci.
\end{enumerate}
\end{matanbox}

\textbf{Definisi istinja secara syariat:} Menghilangkan keluarnya najis \textit{mulawits} (najis yang mengotori yang sifatnya tidak padat) dari \textit{farji} dengan air atau batu. Dan penggunaan air untuk beristinja’ lebih utama.

Istijmar pengertiannya sama dengan istinja.

Makna farji disini adalah:
\begin{enumerate}
    \item Qubul (kemaluan).
    \item Dubur.
\end{enumerate}

Sesuatu yang keluar dari qubul wanita yaitu: 
\begin{enumerate}
    \item Kencing.
    \item Mani.
    \item Madzi. 
    \item Wadzi. 
    \item Darah haid.
    \item Darah istihadhah. 
    \item Darah nifas. 
    \item \textit{ruthubatul farji} (Cairan atau lendir bening/keputihan).
\end{enumerate} 

Sesuatu yang keluar dari dubur yaitu: 
\begin{enumerate}
    \item Tinja.
    \item Kentut. 
    \item Darah. 
    \item Cacing.
\end{enumerate}

Hukum istinja ada 5 yaitu:
\begin{enumerate}
    \item Wajib.
    \item Sunah.
    \item Makruh.
    \item Mubah.
    \item Haram.
\end{enumerate}

Penjelasan hukum istinja:
\begin{enumerate}
    \item Wajib; untuk najis \textit{mulawits} seperti kencing, tinja, madzi, wadzi, dan darah.
    \item Sunah; untuk najis yang padat.
    \item Makruh; untuk kentut.
    \item Mubah untuk keringat dan \textit{rutubah farji} perempuan.
    \item Haram seperti menggunakan air tanpa izin pemiliknya (\textit{ghasab}).
\end{enumerate}